
\typeout{IJCAI--ECAI 26 Instructions for Authors}

% These are the instructions for authors for IJCAI-11.
% They are the same as the ones for IJCAI-07 with superficical wording
%   changes only.

\documentclass{article}
% The file ijcai11.sty is the style file for IJCAI-11 (same as ijcai07.sty).
\usepackage{ijcai26}

% Gregory Mounié: WITH PDFLATEX, original style
% Use the postscript times font!
\usepackage{times}
% Grégory Mounié: WITH LUALATEX INSTEAD
% \usepackage{luatextra}
% \setmainfont{TeX Gyre Termes}

\usepackage{url}
\usepackage[hidelinks]{hyperref}
\usepackage[utf8]{inputenc}
\usepackage{graphicx}
\usepackage{amsmath}
\usepackage{booktabs}
\usepackage[switch]{lineno}

% \usepackage{soul}
% \usepackage[small]{caption}
% \usepackage{amsthm}
%\usepackage{algorithm}
%\usepackage{algorithmic}


% Comment out this line in the camera-ready submission
%\linenumbers

\urlstyle{same}

% the following package is optional:
%\usepackage{latexsym} 

% Following comment is from ijcai97-submit.tex:
% The preparation of these files was supported by Schlumberger Palo Alto
% Research, AT\&T Bell Laboratories, and Morgan Kaufmann Publishers.
% Shirley Jowell, of Morgan Kaufmann Publishers, and Peter F.
% Patel-Schneider, of AT\&T Bell Laboratories collaborated on their
% preparation.

% These instructions can be modified and used in other conferences as long
% as credit to the authors and supporting agencies is retained, this notice
% is not changed, and further modification or reuse is not restricted.
% Neither Shirley Jowell nor Peter F. Patel-Schneider can be listed as
% contacts for providing assistance without their prior permission.

% To use for other conferences, change references to files and the
% conference appropriate and use other authors, contacts, publishers, and
% organizations.
% Also change the deadline and address for returning papers and the length and
% page charge instructions.
% Put where the files are available in the appropriate places.

% PDF Info Is REQUIRED.

% Please leave this \pdfinfo block untouched both for the submission and
% Camera Ready Copy. Do not include Title and Author information in the pdfinfo section
% Grégory: do not work if compile with lualatex
% \pdfinfo{
% /TemplateVersion (IJCAI.2026.0)
% }

\title{Modified IJCAI--ECAI 26 Formatting Instructions for M1Mosig}

% Single author syntax
\author{
    Author Name
    \affiliations
    Affiliation
    \emails
    email@example.com\\
    Supervised by: your laboratory supervisor name
}


\begin{document}




\maketitle

% {% Mosig student
%   {\hbox to0pt{\vbox{\baselineskip=10dd\hrule\hbox
% to\hsize{\vrule\kern3pt\vbox{\kern3pt
% \hbox{{\small I understand what plagiarism entails and I declare that this report }}
% \hbox{{\small is my own, original work. }}
% \hbox{{\small Name, date and signature:}}
% \kern3pt
% }\hfil%\kern3pt
% \vrule
% }\hrule}
% }}
% }

{% Mosig student
\fbox{\parbox{0.9\linewidth}{\small I understand what plagiarism entails and I declare that this report is my own, original work.\\   Name, date and signature:\\ \vspace{0.5cm} }}
}



\begin{abstract}
   This paper provides the style instructions for the M1 Mosig intermediate and final report. \textbf{Do not forgot to include and sign the statement regarding plagiarism.}
\end{abstract}

\section{Introduction}

An introduction


\section{Style and Format}

\LaTeX{} style files that implement these instructions are provided in this archive. (See Section~\ref{stylefiles} for
information about these files.)

\subsection{Layout}

Print manuscripts two columns to a page, in the manner in which these
instructions are printed. The exact dimensions for pages are:
\begin{itemize}
    \item left and right margins: .75$''$
    \item column width: 3.375$''$
    \item gap between columns: .25$''$
    \item top margin---first page: 1.375$''$
    \item top margin---other pages: .75$''$
    \item bottom margin: 1.25$''$
    \item column height---first page: 6.625$''$
    \item column height---other pages: 9$''$
\end{itemize}

All measurements assume an 8-1/2$''$ $\times$ 11$''$ page size. For
A4-size paper, use the given top and left margins, column width,
height, and gap, and modify the bottom and right margins as necessary.


\subsection{Title and Author Information}

Center the title on the entire width of the page in a 14-point bold
font. Below it, center the author name(s) in a 12-point bold font, and
then center the address(es) in a 12-point regular font. Credit to a
sponsoring agency can appear on the first page as a footnote.


\subsection{Abstract}

Place the abstract at the beginning of the first column 3$''$ from the
top of the page, unless that does not leave enough room for the title
and author information. Use a slightly smaller width than in the body
of the paper. Head the abstract with ``Abstract'' centered above the
body of the abstract in a 12-point bold font. The body of the abstract
should be in the same font as the body of the paper.

The abstract should be a concise, one-paragraph summary describing the
general thesis and conclusion of your paper. A reader should be able
to learn the purpose of the paper and the reason for its importance
from the abstract. The abstract should be no more than 200 words long.

\subsection{Text}

The main body of the text immediately follows the abstract. Use
10-point type in a clear, readable font with 1-point leading (10 on
11).

Indent when starting a new paragraph, except after major headings.

\subsection{Headings and Sections}

When necessary, headings should be used to separate major sections of
your paper. (These instructions use many headings to demonstrate their
appearance; your paper should have fewer headings.). All headings should be capitalized using Title Case.

\subsubsection{Section Headings}

Print section headings in 12-point bold type in the style shown in
these instructions. Leave a blank space of approximately 10 points
above and 4 points below section headings.  Number sections with
Arabic numerals.

\subsubsection{Subsection Headings}

Print subsection headings in 11-point bold type. Leave a blank space
of approximately 8 points above and 3 points below subsection
headings. Number subsections with the section number and the
subsection number (in Arabic numerals) separated by a
period.

\subsubsection{Subsubsection Headings}

Print subsubsection headings in 10-point bold type. Leave a blank
space of approximately 6 points above subsubsection headings. Do not
number subsubsections.

\paragraph{Titled paragraphs.} You should use titled paragraphs if and
only if the title covers exactly one paragraph. Such paragraphs should be
separated from the preceding content by at least 3pt, and no more than
6pt. The title should be in 10pt bold font and to end with a period.
After that, a 1em horizontal space should follow the title before
the paragraph's text.

In \LaTeX{} titled paragraphs should be typeset using
\begin{quote}
    {\tt \textbackslash{}paragraph\{Title.\} text} .
\end{quote}

\subsection{Special Sections}

\subsubsection{Appendices}
You may move some of the contents of the paper into one or more appendices that appear after the main content, but before references. These appendices count towards the page limit and are distinct from the supplementary material that can be submitted separately through CMT. Such appendices are useful if you would like to include highly technical material (such as a lengthy calculation) that will disrupt the flow of the paper. They can be included both in papers submitted for review and in camera-ready versions; in the latter case, they will be included in the proceedings (whereas the supplementary materials will not be included in the proceedings).
Appendices are optional. Appendices must appear after the main content.
Appendix sections must use letters instead of Arabic numerals. In \LaTeX, you can use the {\tt \textbackslash{}appendix} command to achieve this followed by  {\tt \textbackslash section\{Appendix\}} for your appendix sections.

%% \subsubsection{Ethical Statement}

%% Ethical Statement is optional. You may include an Ethical Statement to discuss  the ethical aspects and implications of your research. The section should be titled \emph{Ethical Statement} and be typeset like any regular section but without being numbered. This section may be placed on the References pages.

%% Use
%% \begin{quote}
%%     {\tt \textbackslash{}section*\{Ethical Statement\}}
%% \end{quote}

\subsubsection{Acknowledgements}

Acknowledgements are optional.

%% In the camera-ready version you may include an unnumbered acknowledgments section, including acknowledgments of help from colleagues, financial support, and permission to publish. This is not allowed in the anonymous submission. If present, acknowledgements must be in a dedicated, unnumbered section appearing after all regular sections but before references.  This section may be placed on the References pages.

Use
\begin{quote}
    {\tt \textbackslash{}section*\{Acknowledgements\}}
\end{quote}
to typeset the acknowledgements section in \LaTeX{}.


%% \subsubsection{Contribution Statement}

%% Contribution Statement is optional. In the camera-ready version you may include an unnumbered Contribution Statement section, explicitly describing the contribution of each of the co-authors to the paper. This is not allowed in the anonymous submission. If present, Contribution Statement must be in a dedicated, unnumbered section appearing after all regular sections but before references.  This section may be placed on the References pages.

%% Use
%% \begin{quote}
%%     {\tt \textbackslash{}section*\{Contribution Statement\}}
%% \end{quote}
%% to typeset the Contribution Statement section in \LaTeX{}.

\subsubsection{References}

The references section is headed ``References'', printed in the same
style as a section heading but without a number. A sample list of
references is given at the end of these instructions. Use a consistent
format for references, such as that provided by Bib\TeX{}. The reference
list should not include unpublished work.

\subsection{Citations}

Citations within the text should include the author's last name and
the year of publication, for example~\cite{gottlob:nonmon}.  Append
lowercase letters to the year in cases of ambiguity.  Treat multiple
authors as in the following examples:~\cite{abelson-et-al:scheme}
or~\cite{bgf:Lixto} (for more than two authors) and
\cite{brachman-schmolze:kl-one} (for two authors).  If the author
portion of a citation is obvious, omit it, e.g.,
Nebel~\shortcite{nebel:jair-2000}.  Collapse multiple citations as
follows:~\cite{gls:hypertrees,levesque:functional-foundations}.
\nocite{abelson-et-al:scheme}
\nocite{bgf:Lixto}
\nocite{brachman-schmolze:kl-one}
\nocite{gottlob:nonmon}
\nocite{gls:hypertrees}
\nocite{levesque:functional-foundations}
\nocite{levesque:belief}
\nocite{nebel:jair-2000}

\subsection{Footnotes}

Place footnotes at the bottom of the page in a 9-point font.  Refer to
them with superscript numbers.\footnote{This is how your footnotes
    should appear.} Separate them from the text by a short
line.\footnote{Note the line separating these footnotes from the
    text.} Avoid footnotes as much as possible; they interrupt the flow of
the text.

\section{Illustrations}

Place all illustrations (figures, drawings, tables, and photographs)
throughout the paper at the places where they are first discussed,
rather than at the end of the paper.

They should be floated to the top (preferred) or bottom of the page,
unless they are an integral part
of your narrative flow. When placed at the bottom or top of
a page, illustrations may run across both columns, but not when they
appear inline.

Illustrations must be rendered electronically or scanned and placed
directly in your document. They should be cropped outside \LaTeX{},
otherwise portions of the image could reappear during the post-processing of your paper.
When possible, generate your illustrations in a vector format.
When using bitmaps, please use 300dpi resolution at least.
All illustrations should be understandable when printed in black and
white, albeit you can use colors to enhance them. Line weights should
be 1/2-point or thicker. Avoid screens and superimposing type on
patterns, as these effects may not reproduce well.

Number illustrations sequentially. Use references of the following
form: Figure 1, Table 2, etc. Place illustration numbers and captions
under illustrations. Leave a margin of 1/4-inch around the area
covered by the illustration and caption.  Use 9-point type for
captions, labels, and other text in illustrations. Captions should always appear below the illustration.

\section{Tables}

Tables are treated as illustrations containing data. Therefore, they should also appear floated to the top (preferably) or bottom of the page, and with the captions below them.

\begin{table}
    \centering
    \begin{tabular}{lll}
        \hline
        Scenario  & $\delta$ & Runtime \\
        \hline
        Paris     & 0.1s     & 13.65ms \\
        Paris     & 0.2s     & 0.01ms  \\
        New York  & 0.1s     & 92.50ms \\
        Singapore & 0.1s     & 33.33ms \\
        Singapore & 0.2s     & 23.01ms \\
        \hline
    \end{tabular}
    \caption{Latex default table}
    \label{tab:plain}
\end{table}

\begin{table}
    \centering
    \begin{tabular}{lrr}
        \toprule
        Scenario  & $\delta$ (s) & Runtime (ms) \\
        \midrule
        Paris     & 0.1          & 13.65        \\
                  & 0.2          & 0.01         \\
        New York  & 0.1          & 92.50        \\
        Singapore & 0.1          & 33.33        \\
                  & 0.2          & 23.01        \\
        \bottomrule
    \end{tabular}
    \caption{Booktabs table}
    \label{tab:booktabs}
\end{table}

If you are using \LaTeX, you should use the {\tt booktabs} package, because it produces tables that are better than the standard ones. Compare Tables~\ref{tab:plain} and~\ref{tab:booktabs}. The latter is clearly more readable for three reasons:

\begin{enumerate}
    \item The styling is better thanks to using the {\tt booktabs} rulers instead of the default ones.
    \item Numeric columns are right-aligned, making it easier to compare the numbers. Make sure to also right-align the corresponding headers, and to use the same precision for all numbers.
    \item We avoid unnecessary repetition, both between lines (no need to repeat the scenario name in this case) as well as in the content (units can be shown in the column header).
\end{enumerate}

\section{Formulas}

IJCAI's two-column format makes it difficult to typeset long formulas. A usual temptation is to reduce the size of the formula by using the {\tt small} or {\tt tiny} sizes. This doesn't work correctly with the current \LaTeX{} versions, breaking the line spacing of the preceding paragraphs and title, as well as the equation number sizes. The following equation demonstrates the effects (notice that this entire paragraph looks badly formatted, and the line numbers no longer match the text):
%
\begin{tiny}
    \begin{equation}
        x = \prod_{i=1}^n \sum_{j=1}^n j_i + \prod_{i=1}^n \sum_{j=1}^n i_j + \prod_{i=1}^n \sum_{j=1}^n j_i + \prod_{i=1}^n \sum_{j=1}^n i_j + \prod_{i=1}^n \sum_{j=1}^n j_i
    \end{equation}
\end{tiny}%

Reducing formula sizes this way is strictly forbidden. We {\bf strongly} recommend authors to split formulas in multiple lines when they don't fit in a single line. This is the easiest approach to typeset those formulas and provides the most readable output%
%
\begin{align}
    x = & \prod_{i=1}^n \sum_{j=1}^n j_i + \prod_{i=1}^n \sum_{j=1}^n i_j + \prod_{i=1}^n \sum_{j=1}^n j_i + \prod_{i=1}^n \sum_{j=1}^n i_j + \nonumber \\
    +   & \prod_{i=1}^n \sum_{j=1}^n j_i.
\end{align}%

If a line is just slightly longer than the column width, you may use the {\tt resizebox} environment on that equation. The result looks better and doesn't interfere with the paragraph's line spacing: %
\begin{equation}
    \resizebox{.91\linewidth}{!}{$
            \displaystyle
            x = \prod_{i=1}^n \sum_{j=1}^n j_i + \prod_{i=1}^n \sum_{j=1}^n i_j + \prod_{i=1}^n \sum_{j=1}^n j_i + \prod_{i=1}^n \sum_{j=1}^n i_j + \prod_{i=1}^n \sum_{j=1}^n j_i
        $}.
\end{equation}%

This last solution may have to be adapted if you use different equation environments, but it can generally be made to work. Please notice that in any case:

\section{\LaTeX{} Files}\label{stylefiles}

The \LaTeX{} style files are available in this archive.
These style files implement the formatting instructions in this
document.

The \LaTeX{} files are {\tt ijcai26.sty} and {\tt ijcai26.tex}, and
the Bib\TeX{} files are {\tt named.bst} and {\tt ijcai26.bib}. The
\LaTeX{} style file is for version 2e of \LaTeX{}, and the Bib\TeX{}
style file is for version 0.99c of Bib\TeX{} ({\em not} version
0.98i). 

%% The Microsoft Word style file consists of a single file, {\tt
%% ijcai26.docx}. This template is the same as the one used for
%% IJCAI--26.

%% These Microsoft Word and \LaTeX{} files contain the source of the
%% present document and may serve as a formatting sample.

%% Further information on using these styles for the preparation of
%% papers for IJCAI--ECAI 26 can be obtained by contacting {\tt
%%         proceedings@ijcai.org}.

\appendix

%% \section*{Ethical Statement}

%% There are no ethical issues.

\section*{Acknowledgments}

The preparation of these instructions and the \LaTeX{} and Bib\TeX{}
files that implement them was supported by Schlumberger Palo Alto
Research, AT\&T Bell Laboratories, and Morgan Kaufmann Publishers.
Preparation of the Microsoft Word file was supported by IJCAI.  An
early version of this document was created by Shirley Jowell and Peter
F. Patel-Schneider.  It was subsequently modified by Jennifer
Ballentine, Thomas Dean, Bernhard Nebel, Daniel Pagenstecher,
Kurt Steinkraus, Toby Walsh, Carles Sierra, Marc Pujol-Gonzalez,
Francisco Cruz-Mencia and Edith Elkind.


%% The file named.bst is a bibliography style file for BibTeX 0.99c
\bibliographystyle{named}
\bibliography{ijcai26}

\end{document}

